\section*{Rebuttal text}
We thank reviewers for their constructive feedback. We address the common-concerns followed by the individual-concerns. 

Common Concerns:

\textbf{Novelty(Rev-B,C,D,E):}
Prior works target primarily server-class hardware, while MaestoRAG focuses on orchestration of RAG pipeline on edge devices. We observe that replicating prior solutions on edge is insufficient. 
We investigate RAG pipeline through the lens of edge devices and propose a coordinated orchestration strategy. As acknowledged by several reviewers, this is  a timely topic as RAGs are becoming pervasive in edge devices.  
Our fine-grained profiling of RAG-pipeline exposes an optimal-flow(E,RA,G). It also reveals a structural-hazard for pipelining: encoding and generation compete for the GPU. Further, though ‘E’ is compute-bound and given that batch sizes are small, multiple CPU-cores can be deployed for efficient deployment. ‘RA’ is memory-bound and leaves CPU under-utilized. Off-loading ‘RA’ and ‘E’-stages to modern vectorized CPUs removes the hazard and increases utilization. Our adaptive analytical model augments the characterizatization..
Our edge-specific optimizations (aggressive caching, memap’ed indices, shared encoder) reduces working-set size, speeds retrieval, and exploits temporal-locality, yielding gains the baseline doesn’t achieve.

\textbf{Orchestration:
Why encoding-on-CPU?(Rev-B,C,E):} On typical tight-power edge SoCs, the GPU performs encoding and generation, and idles during retrieval. If encoding monopolizes it, subsequent batches must wait for both encoding and generation to finish, minimizing stage-overlap. Worse, keeping separate encoder and generator models resident causes repetitive model swaps within and across batches increases overheads.

\textbf{Throughput-critical vs Latency-Critical(Rev-C):} We note that fixed CPU budget forces a choice: scatter cores across many micro-workers or pack them into a single heavy worker per stage(Section-4.2). Profiling shows extra cores quickly hit diminishing returns for encoding and retrieval, so latency-critical runs use fewer, fatter workers, sacrificing duplicates to shave last milliseconds. 

\textbf{Query-length effect(Rev-B):} With all other hyper-parameters frozen, encoding latency rises just 0.6\% from 10 to 20 tokens and only 1.2\% at 100 tokens; generation is unchanged at 20 tokens and increases 12\% at 100 tokens. (The total input to generation stage is top-5 passages + input query. So system is still getting 500+ tokens.)

\textbf{Scalability(Rev-C,D,E):} We propose a lightweight,  profiling that extracts fine-grained insights and best practices(Section:3-4). Coupled with MaestroRAG’s simple analytical model, the framework can scale to any hardware.

\textbf{Indexing(Rev-A,E):}
\st{We used flat-indexing for accuracy and benchmarked three popular alternatives—IVF-Flat, IVF-PQ, and HNSW—observing performance gains of 29.27\%, 28.15\%, and 32.18\%, respectively, wrt flat-indexing.}
We assume a fixed embedding corpus. Full re-indexing is pushed to low-traffic windows. For new documents, we create delta-indices and query them in parallel. MaestroRAG auto-scales to corpus size, and workload; so any latency is tightly bound. Even mid-merge, end-to-end latency beats PipeRAG’s single-worker setup, and other baselines(Sec:5.1).

\textbf{Platform and Application(Rev-A,C,D):}
Edge-definition: Edge is any battery-powered, user-owned device: phones, tablets, laptops. We model these with Jetson-Orin and low-power desktop PCs. MaestroRAG is built for discrete CPU-GPU machines—zero-copy transfers, aggressive resource allocation, caching, mmap’ed indices—then ported to Jetson’s unified memory, where the optimizations still hold, but gains diminish. Yet it validates the framework’s breadth.
Why wiki-corpus? We target text-only retrieval; combining Wiki-Corpus with the Azure trace gives a  diverse workload that dwarfs on-device data and stress-tests the system. A multimodal, app-centric corpus is future work, but Wiki-Corpus+Azure-trace are sufficient for our use-case.


\textbf{Evaluation:
Accuracy and Caching(Rev-A,B,C):} Supported by many prior works( https://arxiv.org/pdf/2505.11271, https://arxiv.org/pdf/2503.05530 https://arxiv.org/pdf/2502.03771v3, https://arxiv.org/pdf/2406.00025), we enforce a 0.9 semantic-similarity threshold. This retains the factual correctness of the query along with imparting acceleration. Retrieval relies on exact flat-indexing, maximizing recall. Each cache entry carries a user-configurable time-to-live. Both the similarity-threshold and TTL can be tuned.

\textbf{Power(Rev-A):} We conducted power measurements and found MaestroRAG to be more efficient than FlashRAG in terms of average power, peak power and total energy consumed. MaestroRAG consumes 253.3J/query vs 791.8J(+212.6\%) of FlashRAG. We also observed FlashRAG has 16.45\% more peak CPU-power.

\textbf{Ablation(Rev-B,E):}  We observe: Communication latency: 3.36s; No optimization:15.12s, +Adaptive-batching and core allocation:11.79s, +Software optimizations: 6.497s, +Caching: 2.003s.
Cold-start costs(Rev-B): MaestroRAG’s cold-start is ~7s: 0.36s for encoder-load, 3.04s to build vector-indices, and 3.36s to move the generation-model to the GPU. This overhead is one-time before the first batch; FlashRAG re-incurs cold-start with every batch.

\textbf{Individual-Concerns:}

\textbf{Refining writing and plots(Rev-C-E):} Thanks for the specific pointers. We will address said concerns in the camera-ready version if the paper is accepted. 

\textbf{Security implications(Rev-A):} MaestroRAG is agnostic to storage back-ends. For highly sensitive domains, we foresee adopting Differential Privacy-noise or quantised codes to make embeddings themselves non-invertible. But the pipeline orchestration would remain the same.

\textbf{Applicability to other tasks(Rev-A):} Because every stage exchanges only dense embeddings, MaestroRAG is intrinsically task- and language-agnostic. To support retrieval-augmented classification, one simply replaces the generative decoder with a lightweight classifier; for extractive QA, the top-k passages are routed to a span-extraction model such as BERT-QA. In both cases the encoding, retrieval, caching, and scheduling subsystems remain unchanged—and even gain additional headroom. Multilingual deployment is equally straightforward: substitute a cross-lingual encoder (LaBSE or mContriever) and a multilingual generator (mT5 or XGLM). The profiler then re-estimates throughput, and the scheduler  re-allocates CPU/GPU resources, obviating any manual retuning.

\textbf{Memory-constrained environment(Rev-C):} Cache-size is an adjustable hyper-parameter: environments with ample DRAM and high request reuse can provision a large cache, whereas memory-constrained or low-reuse deployments may disable caching entirely. Because SoTA models and their vector stores can saturate main memory, we employ memory-mapped indices to shift most index data to secondary storage, freeing DRAM that can be reassigned for caching. 

\textbf{SOTA Comparison(Rev-D):} While applying elementary optimizations (e.g., fine-grained pipelining) to SoTA systems would be informative, our study purposefully benchmarks each model in its released form to expose the absence of such optimizations. Not every baseline is identical: PipeRAG offers a coarse-grained pipeline that under-utilizes hardware, whereas EdgeRAG builds approximate indices on demand. These contrasts demonstrate that even a well-targeted optimization can outperform more intricate architectures.